\documentclass[10pt,letterpaper]{article}
\usepackage{cornell-notes}

\title{Clause 06.04: Identification Of Concerns}
\author{}
\date{}
\setDocTitle{Clause 06.04: Identification Of Concerns}
\setDocSubtitle{ISO/IEC/IEEE 42010:2022}
\setDocOwner{ISO 42010 Cornell Notes}
\setDocDate{\today}
\setCornellCollection{ISO/IEC/IEEE 42010:2022 Cornell Notes}
\setCornellUnitType{Clause}
\setCornellUnitNumber{06.04}
\setCornellUnitTitle{Identification Of Concerns}
\hypersetup{pdftitle={Clause 06.04: Identification Of Concerns},pdfsubject={Cornell notes},pdfauthor={}}

\begin{document}
\maketitle
\begin{CornellOverviewBox}{Overview}
{\Large\bfseries\textcolor{CNavy}{Section 6.4: Identification of concerns}}\\[3pt]
{\small\textbf{Classification:} Normative}\\
{\small\textbf{Learning objective:} Identify the mandatory, recommended, and permitted provisions governing identification of concerns.}
\end{CornellOverviewBox}
\vspace{3pt}

\begin{CornellNotesTable}
\CornellNoteRow{Purpose}{Establishes how relevant concerns are identified and associated with the stakeholders who hold them.}
\CornellNoteRow{Requirement}{\begin{itemize}[leftmargin=*,nosep,topsep=1pt]
\item A conforming architecture description must identify the concerns considered relevant to the architecture of the entity of interest and consistent with the purpose of the AD.
\item An AD must associate each identified concern with the identified stakeholders holding that concern.
\end{itemize}}
\CornellNoteRow{Example / application}{Concerns include the following: suitability of the architecture for achieving the objectives for the entity of interest, enterprise capability and capacity to implement the entity of interest, the feasibility of realizing and operating the entity of interest, the potential risks and impacts of the entity of interest to its stakeholders throughout its life cycle, added value to the stakeholder(s), reuse of known architectures, resilience, extensibility, adaptability, latency, resource utilization, ...}
\CornellNoteRow{Interpretive note}{\begin{itemize}[leftmargin=*,nosep,topsep=1pt]
\item In general, the association of concerns with stakeholders is many-to-many.
\item Consideration of past, present or future concerns can be relevant to the entity of interest.
\end{itemize}}
\CornellNoteRow{Conformance check}{\begin{itemize}[leftmargin=*,nosep,topsep=1pt]
\item A conforming architecture description must identify the concerns considered relevant to the architecture of the entity of interest and consistent with the purpose of the AD
\item An AD must associate each identified concern with the identified stakeholders holding that concern
\end{itemize}}
\CornellNoteRow{Active recall}{\begin{itemize}[leftmargin=*,nosep,topsep=1pt]
\item Which provisions in Identification of concerns are mandatory for conformance?
\end{itemize}}
\end{CornellNotesTable}


\begin{CornellSummaryBox}{Summary}
Establishes how relevant concerns are identified and associated with the stakeholders who hold them. For conformance, its mandatory provisions must be satisfied, including: A conforming architecture description must identify the concerns considered relevant to the architecture of the entity of interest and consistent with the purpose of the AD; An AD must associate each identified concern with the identified stakeholders holding that concern.
\end{CornellSummaryBox}

\end{document}
